% =====================================================================
% Part II: Output secrecy and the full bound.
%
% Continues the Part I derivation across Alice's decrypt-receive to expose the
% scalar-product output S, builds the absolute-Pr guessing layer on the DERIVED
% endpoint game, transports the Infotheo fiber 1/m to that game through one
% probability identity, and composes with the Part I 2*epsilon_cpa.
%
% Colour / status convention (same as content.tex):
%   \rocqok + \rocq{...}  green : ALREADY EXISTS, cited as a black box.
%   no \rocqok, no \rocq  blue  : TO CONSTRUCT (this plan).
%
% Statement bodies are terse mathematical statements (no proof, no status, no
% motivation).  Proof strategy and the exact Coq names to construct live in
% "% proof:" / "% to create:" comments for the implementer.  Identifiers are
% written in clean math (no \texttt).
%
% Existing module prefixes cited as green:
%   infotheo.dumas2017dual.dsdp.dsdp_game_symbolic   (derived view, Part I)
%   infotheo.dumas2017dual.dsdp.dsdp_indcpa_security (derived 2*epsilon_cpa)
%   infotheo.dumas2017dual.dsdp.dsdp_entropy         (fiber facts)
%   infotheo.dumas2017dual.entropy_fiber.entropy_fiber_zpq
% New identifiers land in the derived chain (dsdp_game_symbolic / dsdp_game_code
% / dsdp_indcpa_security) and a new module dsdp_security_indcpa_fiber.
%
% New notation local to Part II.
\newcommand{\fib}{\mathrm{fiber}}
\newcommand{\guess}{\mathrm{guess}}
\newcommand{\cond}{\mathrm{cond}}
\newcommand{\Prexp}{\Pr\nolimits_{\mathrm{exp}}}
% =====================================================================

\chapter{Extending the derivation to Alice's decryption}
\label{ch:deriving_output}

In the protocol Alice receives the masked aggregate, decrypts it, and removes
her masks to read the scalar-product output $S=U_1V_1+U_2V_2+U_3V_3$. The Part~I
walk stops one reception short of this decrypt
(Definition~\ref{def:walk_obs}), so the derived view exposes only the
ciphertexts. This chapter continues the same walk across that decrypt-reception,
lowers and denotes the new observation, and so the derived endpoint game exposes
$S$ alongside the ciphertexts. Adjoining $S$ to both endpoints is a common
deterministic post-processing: it adds no encryption hop, so the Part~I bound
$2\,\epscpa$ carries over unchanged.

\begin{definition}[Output reception]
  \label{def:alice_obs_output}
  \uses{def:alice_obs}
  % to create: alice_obs constructor AO_recv_output (dsdp_game_symbolic).
  The corrupted-view vocabulary (Definition~\ref{def:alice_obs}) gains one
  action: receive the decrypted output and bind it to a result name. It is not
  an encryption hop.
  \[ o \mathrel{::=} \dots \mid \mathsf{recvOut}(x) \]
\end{definition}

\begin{definition}[Walking across the decrypt-reception]
  \label{def:walk_obs_output}
  \uses{def:walk_obs,def:alice_obs_output}
  % to create: extend walk_obs (dsdp_game_symbolic) with the decrypt-receive arm
  %   and feed the three-element response stream dsdp_recv_responses (the third
  %   entry HE_var 50 is Alice's decrypt input); emit AO_recv_output carrying the
  %   he_term for S, then halt.
  At the decrypt-reception the extended walk reads the received plaintext off the
  response stream and emits the output reception instead of halting. The
  ciphertext receptions and assemblies before it are unchanged, so the
  encryption-hop count is unchanged.
\end{definition}

\begin{lemma}[The derived output-exposing DSDP trace]
  \label{lem:obs_of_procs_dsdp_leak_S}
  \uses{lem:obs_of_procs_dsdp,def:walk_obs_output}
  % to create: obs_of_procs_dsdp_leak_S, the equational characterisation of the
  %   extended obs_of_procs on the DSDP inputs; computes to the Part I trace plus
  %   one AO_recv_output before the leak.
  For the DSDP inputs the extended trace computes to the Part~I trace
  (Lemma~\ref{lem:obs_of_procs_dsdp}) with the output reception inserted before
  the leak. The two encryption hops are exactly the Part~I two.
  \[ \tau^{S}_{\mathrm{DSDP}}
        = \tau_{\mathrm{DSDP}} \text{ with } \mathsf{recvOut}(S)
        \text{ before } \mathsf{leak} \]
\end{lemma}

\begin{definition}[Lowering the output reception]
  \label{def:lower_obs_output}
  \uses{def:lower_obs,def:alice_obs_output}
  % to create: extend lower_obs (dsdp_game_symbolic) with the AO_recv_output arm
  %   -> a let binding the resolved S term followed by a put to S_output_cell;
  %   reuses existing game_code constructors (let, put), so game_code is NOT
  %   extended.
  The lowering pass maps the output reception to a let-binding of the resolved
  output term followed by a write to a dedicated output cell. It introduces no
  new game-code statement and no new hop.
  \[ \mathsf{recvOut}(S) \;\mapsto\; \mathsf{let}(S);\ \mathsf{put}_{S}(S) \]
\end{definition}

\begin{definition}[Denoting the output and exposing it]
  \label{def:denote_game_leak_S}
  \uses{def:denote_game,def:lower_obs_output}
  % to create: S_output_cell : Location; id_s_get oracle id; denote_s_get_body;
  %   widen the denoted game's export interface with #val #[id_s_get] : 'unit ->
  %   t_msg returning S_output_cell.  GC_ret stays cipher_list, so the Part I
  %   generic theorem is untouched.
  Denotation writes the lowered output to a dedicated cell and adds one oracle
  reading it, leaving the return value and every other statement as in Part~I
  (Definition~\ref{def:denote_game}).
  \[ \sem{\mathsf{put}_{S}(S)} : \text{write } S \text{ to } S\text{-cell},
        \qquad \mathsf{getOut} : \mathbf{1} \to \mathsf{msg} \]
\end{definition}

\begin{definition}[The output-exposing endpoint games]
  \label{def:games_leak_s}
  \uses{def:denote_game_leak_S,def:games}
  % to create: real_game_leak_S, zero_game_leak_S (dsdp_indcpa_security), the
  %   real/all-zero denotations of the output-exposing DSDP problem.
  The real and all-zero games of the DSDP problem
  (Definition~\ref{def:games}) denoted from the output-exposing derivation. Each
  exposes the ciphertexts and the output $S$. Write $\real_S$ and $\zero_S$ for
  the two.
  \[ \view = \big(\,[\,a_1,a_2,c_2,c_3\,],\ S\,\big),\qquad
        S = u_1 v_1 + u_2 v_2 + u_3 v_3 \]
\end{definition}

\begin{lemma}[Computational distance, output-exposing games]
  \label{lem:dsdp_advantage_derived_leak_S}
  \uses{def:games_leak_s,thm:dsdp_secure}
  % to create: dsdp_advantage_derived_leak_S; the output cell is a common
  %   deterministic post-processing of both endpoints (no new hop), so the
  %   Part I 2*epsilon_cpa bound (thm:dsdp_secure) transports verbatim.
  The advantage between the output-exposing real and all-zero games is at most
  $2\,\epscpa$, the Part~I bound (Theorem~\ref{thm:dsdp_secure}).
  \[ \AdvE(\,\real_S,\ \zero_S,\ A\,) \le 2\,\epscpa \]
\end{lemma}

\chapter{The output-secrecy guessing experiment}
\label{ch:guessing_experiment}

The output channel is measured not by an advantage but by the absolute
probability that a predictor, reading the view including $S$, names Bob's secret
$V_2$. This chapter builds that guessing experiment on the derived endpoint game
and rewrites its success probability as an Infotheo distribution probability, the
single identity that carries the fiber bound of Chapter~\ref{ch:it_bound} to the
\ssprove{} side.

\begin{definition}[Guessing experiment on the endpoint game]
  \label{def:guessing_experiment_leak_S}
  \uses{def:games_leak_s}
  % to create: predictor_guesser, guessing_challenger, guessing_experiment
  %   (dsdp_security_indcpa_fiber).  Challenger imports id_game_run, id_s_get,
  %   id_v2_get; feeds (ciphertexts, S) to the predictor; tests guess == V2.
  A predictor reads the leaked ciphertexts and the output $S$ and returns a
  guess. The experiment runs it against the endpoint game and tests
  $\guess = V_2$. Its success probability is written $\Prexp(\text{game})$.
  \[ \Prexp(\text{game})
        = \mu\big(\mathrm{Pr}(\,\text{challenger}\circ\text{predictor}
        \circ\text{game}\,)\big)(\mathrm{true}) \]
\end{definition}

\begin{lemma}[\ssprove{} absolute-$\mathrm{Pr}$ reflection calculus (library)]
  \label{lem:exist_pr_code}
  \rocqok
  % library: Pr_code_ret/_sample/_bind/_get/_put (Crypt.nominal.Pr), dlet_uniform
  %   (Crypt.HybridArgument), SDistr = {distr A/R}, dletE, dunit1E, Pr_Pr_fst.
  SSProve's nominal layer expresses the probability of a closed program as an
  iterated sum over its sample operations, a uniform draw of domain $n$
  contributing the factor $1/n$. Its subdistribution type is the
  mathcomp-analysis distribution $\{\mathrm{distr}\ A/R\}$.
  \[ \mathrm{Pr}(x \leftarrow \mathrm{sample}\ \mathrm{unif}\,n \,;;\, k\,x)
        = \tfrac1n{\textstyle\sum_x}\mathrm{Pr}(k\,x) \]
\end{lemma}

\begin{definition}[The guessing-experiment distribution]
  \label{def:dsdp_guess_fdist}
  \uses{def:guessing_experiment_leak_S}
  % to create: dsdp_guess_fdist (dsdp_security_indcpa_fiber), the uniform fdist
  %   on the experiment's sample tuple; V2,V3,U2,U3,S,guess as RVs over it.
  The uniform distribution $P$ on the samples of the guessing experiment at
  $\zero_S$, the game's draws together with the predictor's. The protocol
  quantities and the predictor's output are random variables over $P$.
  \[ P : \{\mathrm{fdist}\ \mathrm{Samples}\},
        \qquad V_2,\,V_3,\,U_2,\,U_3,\,S,\,\guess : \{\mathrm{RV}\ P\} \]
\end{definition}

\begin{lemma}[The game probability is the fdist probability]
  \label{lem:Pr_guess_enc_zero_leak_S_eqE}
  \uses{def:games_leak_s,def:dsdp_guess_fdist,def:guessing_experiment_leak_S,lem:exist_pr_code}
  % to create: Pr_guess_enc_zero_leak_S_eqE (the connector).  Proof: the closed
  %   experiment denotes to the uniform distribution on its samples, which is P
  %   (the predictor's draws are extra coordinates of P, so an opaque predictor
  %   needs no inspection); reflection (lem:exist_pr_code) rewrites both sides to
  %   the uniform average of [guess = V2].  ~30-80 lines, one induction over the
  %   sample list.
  $\Prexp(\zero_S)$ equals the Infotheo probability, over $P$, of the event
  $\guess = V_2$. Both reduce to the same uniform average of the indicator
  $[\,\guess = V_2\,]$ over the experiment's samples. This identity carries the
  fiber bound, proved on the Infotheo side, to $\Prexp(\zero_S)$.
  \[ \Prexp(\zero_S) = \Pr\nolimits_{P}[\,\guess = V_2\,] \]
\end{lemma}

\chapter{The information-theoretic bound and full secrecy}
\label{ch:it_bound}

Conditioned on $S$ and her inputs, Bob's secret $V_2$ is uniform over a fiber of
$m$ points, so any guess succeeds with probability at most $1/m$. Infotheo proves
this over an abstract joint distribution. Instantiating that distribution at the
guessing-experiment distribution $P$ and pushing the bound through the connector
(Lemma~\ref{lem:Pr_guess_enc_zero_leak_S_eqE}) gives the output bound at the
endpoint, and composing with the computational distance gives the full secrecy
bound $1/m + 2\,\epscpa$.

\section{Reused Infotheo facts}

\begin{lemma}[Fiber cardinality is $m$]
  \label{lem:exist_fiber_card}
  \rocq{infotheo.dumas2017dual.dsdp.dsdp_entropy.dsdp_fiber_card}
  \rocqok
  The solution set of the scalar-product constraint, for admissible
  multipliers, has exactly $m$ points.
  \[ \#\lvert \fib(u_1,u_2,u_3,v_1,s)\rvert = m \]
\end{lemma}

\begin{lemma}[Conditional uniformity on the fiber]
  \label{lem:exist_cpr_fiber}
  \rocq{infotheo.dumas2017dual.entropy_fiber.entropy_fiber_zpq.cPr_uniform_fiber}
  \rocqok
  \uses{lem:exist_fiber_card}
  For a joint distribution in which the variable pair is uniform and
  independent of the inputs, conditioning on the revealed values leaves the
  variable uniform over its fiber.
  \[ \Pr[\,\mathrm{Var}=v \mid \mathrm{Cond}=\cond\,]
        = \#\lvert\fib(\cond)\rvert^{-1} \]
\end{lemma}

\begin{lemma}[DSDP solution is uniform, probability $1/m$]
  \label{lem:exist_pr_sol}
  \rocq{infotheo.dumas2017dual.dsdp.dsdp_entropy.Pr_dsdp_sol_uniform}
  \rocqok
  \uses{lem:exist_cpr_fiber,lem:exist_fiber_card}
  Each solution $(v_2,v_3)$ of the DSDP constraint has conditional probability
  $1/m$ given the inputs and $S$.
  \[ \Pr[\,(V_2,V_3)=(v_2,v_3)\mid \mathrm{Cond}=\cond\,] = m^{-1} \]
\end{lemma}

\begin{theorem}[Residual entropy is $\log m$]
  \label{thm:exist_centropy}
  \rocq{infotheo.dumas2017dual.dsdp.dsdp_entropy.dsdp_centropy_uniform}
  \rocqok
  \uses{lem:exist_pr_sol}
  Conditioned on the inputs and $S$, the variable pair retains $\log m$ bits of
  entropy.
  \[ H(\mathrm{Var}\mid \mathrm{Cond}) = \log m \]
\end{theorem}

\section{Instantiating the fiber bound at $P$}

\begin{lemma}[Distribution identities: uniformity, independence, constraint]
  \label{lem:dsdp_guess_identities}
  \uses{def:dsdp_guess_fdist,def:games_leak_s}
  % to create: dsdp_guess_V2_uniform, dsdp_guess_VarRV_indep_inputs,
  %   dsdp_guess_S_determined, dsdp_guess_indep_V2_given_S
  %   (dsdp_security_indcpa_fiber); the hypotheses dsdp_entropy needs, read off
  %   the uniform structure of P and the all-zero ciphertexts.
  Over $P$, the variable pair $(V_2,V_3)$ is uniform and independent of the
  inputs, the output $S$ is the constraint function of the variables and inputs,
  and the guess depends on $V_2$ only through $S$.
  \[ {p}_{(V_2,V_3)} = \mathrm{unif},\quad
     (V_2,V_3)\perp \mathrm{Inputs},\quad
     S = g(\mathrm{Var},\mathrm{Input}),\quad
     \guess \perp V_2 \mid S \]
\end{lemma}

\begin{lemma}[Fiber residual]
  \label{lem:Pr_guess_fiber_le_invm}
  \uses{lem:exist_pr_sol,lem:dsdp_guess_identities}
  % to create: Pr_guess_fiber_le_invm (dsdp_security_indcpa_fiber).  Proof:
  %   condition on S; lem:exist_pr_sol gives each V2 value probability 1/m, and
  %   guess _|_ V2 | S (lem:dsdp_guess_identities) caps the collision at 1/m.
  Over $P$, the predictor guesses $V_2$ with probability at most $1/m$.
  \[ \Pr\nolimits_{P}[\,\guess=V_2\,] \le 1/m \]
\end{lemma}

\section{The output bound and composition}

\begin{theorem}[Information-theoretic bound at the endpoint]
  \label{thm:Pr_guess_enc_zero_leak_S_le_invm}
  \uses{lem:Pr_guess_enc_zero_leak_S_eqE,lem:Pr_guess_fiber_le_invm}
  % to create: Pr_guess_enc_zero_leak_S_le_invm (dsdp_security_indcpa_fiber).
  %   Proof: rewrite Prexp(zero_S) by the connector, then the fiber residual.
  At the output-exposing all-zero endpoint, the predictor guesses $V_2$ with
  probability at most $1/m$.
  \[ \Prexp(\zero_S) \le 1/m \]
\end{theorem}

\begin{theorem}[Full secrecy by composition]
  \label{thm:dsdp_alice_secrecy_leak_S}
  \uses{thm:Pr_guess_enc_zero_leak_S_le_invm,lem:dsdp_advantage_derived_leak_S}
  % to create: dsdp_alice_secrecy_leak_S (dsdp_security_indcpa_fiber).  Proof:
  %   triangle Prexp(real_S) <= Prexp(zero_S) + AdvE(real_S, zero_S, predictor),
  %   then thm:Pr_guess_enc_zero_leak_S_le_invm (1/m) and
  %   lem:dsdp_advantage_derived_leak_S (2*epsilon_cpa).
  Against the output-exposing real game, the predictor guesses $V_2$ with
  probability at most $1/m + 2\,\epscpa$.
  \[ \Prexp(\real_S) \le 1/m + 2\,\epscpa \]
\end{theorem}
